\section{Conclusion}
\label{sec:conclusion}

% VALUE: the insights at a high level: the compiled order is the schedule of an array that one in-order CP feeds, so its false dependences can be read, bounded, and removed after compilation; the two clauses, read from the artifacts, say where removal can pay, so the cut decides where and the prompt length how much; OPCORD reads, rewrites behind checkers, and prices order and configuration together per prompt length; on AIR's designs it shortens prefill and decode with bit-identical outputs and explains why a closed engine's flash prefill leaves little; other in-order issuers feeding spatial engines may face the same choice; the speedups are snapshots, the insight and the condition durable, and an improved design is a new input
% DECISION: the transfer to other accelerators is a scope claim, worded as a possibility (step 4 section 2); the closed engine's explanation is claimed for its flash prefill only until the FLM per-operator merge replay lands (step 2, section 5 item 1).
When one in-order command processor feeds all the engines of a spatial NPU, the order compiled into a design is the schedule that the array runs, so the false dependences that leave its tiles idle can be read, bounded, and removed from the compiled design itself.
Whether removing one frees overlap depends on whether the freed and the awaited work use disjoint resources and whether no true dependence separates them, which the artifacts show before any duration is measured, so a design's cut fixes where overlap can be freed and the prompt length fixes what it is worth.
OPCORD reads compiled designs, open or closed, as dependences among the CP, the DMA channels, the cores, and the host, bounds each false dependence, and rewrites issue order and buffers only where its checkers show the rewrite legal.
On AIR's fused designs for five LLMs, these rewrites shorten both prefill and decode with bit-identical outputs, and the same reading explains why the flash prefill of a closed engine's current release leaves little to recover.
Because the order and the configuration that pay both change with the prompt, OPCORD chooses them together for each prompt length with an event model rather than by running every candidate.
Other accelerators in which one in-order issuer feeds spatially partitioned engines may face the same choice.
The speedups and the stack comparisons are snapshots of current releases, whereas the insight and the condition hold as designs improve, and each improved design is a new input to OPCORD rather than a replacement for it.
